\section{Detailed methods (implementation-aligned)}
\label{app:methods}

This appendix records implementation-aligned algorithmic details for \citrees.
It is intended to make the manuscript reproducible and to keep the main text
readable. Proofs for the paper's fixed-node/root claims are given in
Appendices~A--G.

\subsection{Algorithm overview}
\label{app:methods-overview}

\citrees implements conditional inference trees and forests that use
permutation-based hypothesis testing for variable and split selection. The core
design choice relative to CART-style trees is to separate variable screening
(Stage~A) from split-point optimization (Stage~B) and to attach p-values to the
screening step at a fixed node (especially the root).

\subsection{Conditional inference tree (high level)}
\label{app:methods-tree}

\begin{algorithm}[t]
  \caption{Conditional inference tree (Stage A $\rightarrow$ Stage B)}
  \label{alg:citree}
  \begin{algorithmic}[1]
    \Require Node data $(X_t, y_t)$, significance levels $\alpha_{\mathrm{sel}}, \alpha_{\mathrm{split}}$
    \State \textbf{if} stopping criteria hold \textbf{then} \Return leaf value from $y_t$
    \State $F_t \gets$ candidate features (drop node-constant features; apply \texttt{max\_features} if configured)
    \State \textbf{Stage A:} compute $(p_{t,j})_{j \in F_t}$ via +1 permutation tests; set $j_t^\star \gets \argmin_{j \in F_t} p_{t,j}$
    \State \textbf{if} $\min_{j \in F_t} p_{t,j} \ge \alpha_{\mathrm{sel}}/|F_t|$ \textbf{then} \Return leaf
    \State $C_{t,j_t^\star} \gets$ candidate thresholds for feature $j_t^\star$
    \State \textbf{Stage B:} choose $c_t^\star \in C_{t,j_t^\star}$ (optionally with a permutation test on split quality)
    \State \textbf{if} Stage B fails (e.g., not significant or violates \texttt{min\_impurity\_decrease} / min-leaf constraints) \textbf{then} \Return leaf
    \State Split $(X_t, y_t)$ into left/right children using $(j_t^\star, c_t^\star)$
    \State Recurse on children and return an internal node
  \end{algorithmic}
\end{algorithm}

\paragraph{Notes.}
If \texttt{adjust\_alpha\_selector} / \texttt{adjust\_alpha\_splitter} are
disabled, compare to $\alpha_{\mathrm{sel}}$ / $\alpha_{\mathrm{split}}$ rather
than Bonferroni-adjusted levels. If permutation testing is disabled for a stage
(e.g., \texttt{n\_resamples\_selector = None}), that stage uses raw
selector/splitter scores without p-values.

\subsection{Stopping criteria}
\label{app:methods-stopping}

At a node $t$, tree growth stops when any of the following conditions hold:
\begin{enumerate}
  \item $n_t < \texttt{min\_samples\_split}$,
  \item depth exceeds \texttt{max\_depth},
  \item the node is pure (classification) or otherwise terminal per configuration,
  \item Stage~A fails to reject any null hypothesis (no significant features),
  \item Stage~B fails to validate any split (e.g., not significant),
  \item impurity decrease is below \texttt{min\_impurity\_decrease},
  \item a candidate split would violate \texttt{min\_samples\_leaf}.
\end{enumerate}

\subsection{Leaf values}
\label{app:methods-leaves}

\paragraph{Classification.}
For a leaf containing indices $L$, the predicted class probabilities are
%
\begin{equation}
  \hat{p}_k(L) = \frac{1}{|L|} \sum_{i \in L} \ind\{y_i = k\}, \qquad k=1,\dots,K.
\end{equation}

\paragraph{Regression.}
For a leaf containing indices $L$, the leaf prediction is the mean response,
%
\begin{equation}
  \hat{\mu}(L) = \frac{1}{|L|} \sum_{i \in L} y_i.
\end{equation}

\subsection{Stage A: feature screening}
\label{app:methods-stageA}

At node $t$ with sample indices $I_t$ and data $(X_t, y_t)$, Stage~A tests each
candidate feature $X_{t,j}$ for association with $y_t$. Features that are
constant at the node are removed from the candidate set since they cannot
become non-constant within the subtree rooted at $t$.

\subsubsection{Null hypothesis}
\label{app:methods-stageA-null}

For each feature $j \in F_t$,
%
\begin{equation}
  H^{\mathrm{sel}}_{t,j}: X_{t,j} \perp y_t.
\end{equation}

\subsubsection{Selector statistics}
\label{app:methods-selectors}

\paragraph{Multiple correlation (mc; classification).}
Define
%
\begin{equation}
  \mathrm{mc}(x,y) := \sqrt{\frac{\mathrm{SSB}}{\mathrm{SST}}},
\end{equation}
%
where $\mathrm{SST} = \sum_{i=1}^n (x_i - \bar{x})^2$ and
$\mathrm{SSB} = \sum_{k=1}^K n_k(\bar{x}_k - \bar{x})^2$.
This score lies in $[0,1]$ and is a monotone transform of the one-way ANOVA $F$
statistic at fixed $(n,K)$, so fixed-$B$ permutation p-values (with the same tie
convention) agree under either statistic.

\paragraph{Mutual information (mi; classification).}
Mutual information satisfies $I(X;Y)\ge 0$ with $I(X;Y)=0$ iff $X \perp Y$.
\citrees uses a $k$-nearest-neighbor estimator \citep{Kraskov2004EstimatingMutualInformation}.
Because MI is not bounded on $[0,1]$, it is not mixed with bounded selectors in
max-statistic mode without an explicit normalization policy.

\paragraph{Pearson correlation (pc; regression).}
Define
%
\begin{equation}
  \mathrm{pc}(x,y) := \left|\frac{\sum_{i=1}^n (x_i-\bar{x})(y_i-\bar{y})}
  {\sqrt{\sum_{i=1}^n (x_i-\bar{x})^2}\sqrt{\sum_{i=1}^n (y_i-\bar{y})^2}}\right|.
\end{equation}
%
This lies in $[0,1]$ and is a monotone transform of the usual correlation $t$
statistic at fixed $n$, so it yields the same fixed-$B$ permutation p-values
under the same tail/tie conventions.

\paragraph{Distance correlation (dc; regression).}
Distance correlation detects both linear and nonlinear dependence and satisfies
$\mathrm{dCor}(X,Y)=0$ iff $X \perp Y$ under mild moment conditions
\citep{Szekely2007DistanceCorrelation}.
\citrees uses an external implementation of distance correlation.

\paragraph{Randomized dependence coefficient (rdc; classification and regression).}
RDC is a computationally efficient nonlinear dependence score based on random
nonlinear projections of rank-transformed data \citep{LopezPaz2013RDC}. In
\citrees, RDC is used as a bounded score in $[0,1]$ and can be combined with
other bounded selectors in multi-selector mode.

\subsubsection{Permutation p-values}
\label{app:methods-stageA-pvalues}

For each feature $j \in F_t$, Stage~A computes a Monte Carlo permutation p-value
using the Phipson--Smyth (+1) correction as described in \cref{sec:plusone}. The
paper's formal validity claims for these p-values are limited to fixed nodes
and fixed-$B$ testing (see \cref{sec:theory} and Appendices~A--G).

\subsection{Stage B: threshold screening (scope note)}
\label{app:methods-stageB}

Given the selected feature $j_t^\star$, \citrees considers a finite threshold
set $C_{t,j_t^\star}$ and chooses a split point by optimizing split quality
over candidates (and optionally performing a permutation test over thresholds).
Because Stage~B is applied after selecting $j_t^\star$ using the same labels,
we treat Stage~B p-values as post-selection/adaptive statistics unless
additional selective-inference machinery or sample splitting is used.

\subsubsection{Null hypothesis (conceptual)}
\label{app:methods-stageB-null}

For a fixed threshold $c \in C_{t,j_t^\star}$, define the induced partition
$I_t^L(c) := \{i \in I_t : X_{i,j_t^\star} \le c\}$ and
$I_t^R(c) := I_t \setminus I_t^L(c)$. The conceptual null for a split-quality
test at threshold $c$ is that the labels are exchangeable with respect to this
partition (i.e., there is no association between $y_t$ and the left/right child
assignment). In the adaptive tree algorithm, this is tested only after choosing
$j_t^\star$ using the same $y_t$, so the resulting p-values are post-selection.

\subsubsection{Split quality statistics (splitters)}
\label{app:methods-splitters}

\paragraph{Convention (implementation-aligned).}
For permutation testing of split quality, \citrees uses the \emph{unweighted}
sum of child impurities, $I(y_L) + I(y_R)$, with a left-tail convention (smaller
is better). Weighted impurity is used for deterministic threshold scanning and
for the \texttt{min\_impurity\_decrease} constraint.

\paragraph{Classification impurities.}
For class probabilities $\hat p_k = \frac{1}{n}\sum_{i=1}^n \ind\{y_i=k\}$,
%
\begin{equation}
  \mathrm{Gini}(y) = 1 - \sum_{k=1}^K \hat p_k^2,
  \qquad
  \mathrm{Entropy}(y) = -\sum_{k=1}^K \hat p_k \log_2(\hat p_k).
\end{equation}
%
For a threshold $c$ inducing children $(y_L,y_R)$, the tested split statistic is
%
\begin{equation}
  T^{\mathrm{split}}(c) := I(y_L) + I(y_R),
\end{equation}
%
where $I(\cdot)$ is either Gini or Entropy.

\paragraph{Regression impurities.}
For regression,
%
\begin{equation}
  \mathrm{MSE}(y) = \frac{1}{n}\sum_{i=1}^n (y_i-\bar y)^2,
  \qquad
  \mathrm{MAE}(y) = \frac{1}{n}\sum_{i=1}^n \big|y_i - \mathrm{median}(y)\big|.
\end{equation}

\subsubsection{Threshold candidate sets}
\label{app:methods-thresholds}

\citrees generates candidate thresholds from the set of midpoints between
consecutive unique feature values. Let $x_{(1)} < \cdots < x_{(m)}$ denote the
sorted unique values of a feature at node $t$; the full midpoint set is
%
\begin{equation}
  \mathrm{Midpoints}(x) := \left\{\frac{x_{(i)} + x_{(i+1)}}{2} : i=1,\dots,m-1\right\}.
\end{equation}

\paragraph{Threshold methods.}
\citrees supports several threshold selection policies:
\begin{enumerate}
  \item \textbf{Exact (default):} use all midpoints,
  \item \textbf{Random:} subsample midpoints without replacement,
  \item \textbf{Percentile:} choose evenly spaced percentiles of the midpoint set,
  \item \textbf{Histogram:} choose histogram-based bin edges of the midpoint set.
\end{enumerate}

\paragraph{Small-sample fallback.}
If a feature has very few unique values (implementation default: $\le 4$),
\citrees uses all midpoints regardless of the configured threshold method.

\paragraph{\texttt{max\_thresholds} semantics.}
When \texttt{max\_thresholds} is set, the requested count is computed by the
same integer/float/\texttt{sqrt}/\texttt{log2} rules used by \texttt{max\_features},
then applied to the candidate midpoints for the chosen threshold method. When
\texttt{threshold\_method="exact"}, all midpoints are returned regardless of
\texttt{max\_thresholds} (accepted for API compatibility but not used).

\subsubsection{Split selection procedure (pseudocode)}
\label{app:methods-stageB-alg}

\begin{algorithm}[t]
  \caption{Split selection at a node (Stage B)}
  \label{alg:stageB}
  \begin{algorithmic}[1]
    \Require Feature vector $x = X_{t,j_t^\star}$, labels $y_t$, level $\alpha_{\mathrm{split}}$, candidate thresholds $C_{t,j_t^\star}$
    \Require Resamples $B$, \texttt{adjust\_alpha\_splitter}, \texttt{min\_samples\_leaf}, optional early stopping
    \State Generate candidate thresholds $C_{t,j_t^\star}$ from midpoints (filter those violating \texttt{min\_samples\_leaf})
    \If{$|C_{t,j_t^\star}| = 0$}
      \State \Return no split
    \EndIf
    \State $\alpha_{\mathrm{eff}} \gets \alpha_{\mathrm{split}}/|C_{t,j_t^\star}|$ if Bonferroni is enabled else $\alpha_{\mathrm{split}}$
    \State Initialize best threshold $c^\star$ and best p-value $p^\star \gets 1$
    \For{each $c \in C_{t,j_t^\star}$ (optionally sorted by a scanning heuristic)}
      \State Compute observed split score $T_0(c) \gets I(y_L(c)) + I(y_R(c))$
      \State Compute permutation p-value $p_c$ for the left tail using $B$ label permutations
      \If{$p_c < p^\star$}
        \State $(p^\star, c^\star) \gets (p_c, c)$
      \EndIf
      \If{early stopping enabled and $p^\star < \alpha_{\mathrm{eff}}$}
        \State \textbf{break}
      \EndIf
    \EndFor
    \State Apply deterministic constraints (e.g., \texttt{min\_impurity\_decrease}); if violated, return no split
    \State \Return split $(c^\star, p^\star)$
  \end{algorithmic}
\end{algorithm}

\subsection{Permutation testing (fixed-B baseline)}
\label{app:methods-permutation}

For a scalar test statistic $T(\cdot,\cdot)$ with a right-tail convention, let
$T_0$ denote the observed statistic and let $T_1,\dots,T_B$ denote the statistic
computed on $B$ i.i.d.\ random permutations of the labels. \citrees uses the
Phipson--Smyth (+1) Monte Carlo permutation p-value \citep{PhipsonSmyth2010PermutationPValues},
%
\begin{equation}
  p := \frac{1 + \sum_{b=1}^B \ind\{T_b \ge T_0\}}{B+1},
\end{equation}
%
with the inequality direction chosen to match the tail convention. Under the
fixed-node assumptions in \cref{app:assumptions}, these +1 p-values are
super-uniform in fixed-$B$ mode (\cref{sec:theory}).

\paragraph{Resolution in fixed-$B$ mode.}
Because $p$ lies on the grid $\{1/(B+1),2/(B+1),\dots,1\}$, any configured
significance threshold below $1/(B+1)$ is unattainable. In particular, if Stage~A
uses Bonferroni at level $\alpha_{\mathrm{sel}}/m_t$, then a necessary condition
for any rejection is $\alpha_{\mathrm{sel}}/m_t \ge 1/(B+1)$ (equivalently
$B \ge \lceil m_t/\alpha_{\mathrm{sel}} \rceil - 1$). Discreteness can also
create ties (many candidates achieving the same minimal p-value), so a
tie-breaking convention must be specified.

\subsection{Sequential permutation testing (early stopping; scope note)}
\label{app:methods-sequential}

\citrees implements sequential stopping rules to reduce the computational cost
of permutation testing. These are useful engineering heuristics, but they change
the inferential object: if the algorithm stops at a data-dependent time and then
reports the running +1 estimate $(k+1)/(b+1)$, this value should not be treated
as a classical fixed-$B$ Monte Carlo p-value without additional sequential
analysis. For publication-grade p-value claims in this manuscript, use fixed-$B$
mode (\texttt{early\_stopping\_* = None}) and report $B$ explicitly.

\subsubsection{Simple sequential stopping (baseline)}
\label{app:methods-sequential-simple}

The ``simple'' rule stops early either when the running estimate is already
below the relevant threshold (significance) or when rejection is impossible
given the remaining budget (futility). This rule is fast but can inflate Type I
error by repeatedly checking the running estimate without a sequential
adjustment.

\begin{algorithm}[t]
  \caption{Simple sequential permutation test (heuristic)}
  \label{alg:simple-sequential}
  \begin{algorithmic}[1]
    \Require Statistic $T$, observed value $T_0$, max resamples $B$, threshold $\alpha$
    \State $\texttt{min\_resamples} \gets \lceil 1/\alpha \rceil$; $B \gets \max(B,\texttt{min\_resamples})$; $k \gets 0$
    \For{$b = 1,\dots,B$}
      \State Sample a random permutation and compute $T_b$
      \If{$T_b \ge T_0$} \State $k \gets k+1$ \EndIf
      \State $p_{\mathrm{cur}} \gets (k+1)/(b+1)$
      \If{$b \ge \texttt{min\_resamples}$}
        \If{$p_{\mathrm{cur}} < \alpha$} \State \Return $p_{\mathrm{cur}}$ \EndIf
        \State $p_{\mathrm{best}} \gets (k+1)/(B+1)$
        \If{$p_{\mathrm{best}} \ge \alpha$} \State \Return $p_{\mathrm{cur}}$ \EndIf
      \EndIf
    \EndFor
    \State \Return $(k+1)/(B+1)$
  \end{algorithmic}
\end{algorithm}

\paragraph{Left-tail variant (Stage B).}
For split testing, the same template is used with a left-tail convention, i.e.,
replace the exceedance check by $T_b \le T_0$.

\subsubsection{Adaptive sequential stopping (Bayesian)}
\label{app:methods-sequential-adaptive}

The ``adaptive'' rule uses a Beta--Binomial model to stop once the algorithm is
confident that the (latent) p-value is either below or above the threshold. Let
$k$ be the number of exceedances after $b$ permutations (right tail), model
%
\begin{equation}
  k \mid p \sim \mathrm{Binomial}(b,p), \qquad p \sim \mathrm{Beta}(1,1),
  \qquad p \mid k,b \sim \mathrm{Beta}(1+k, 1+b-k).
\end{equation}
%
Given a confidence parameter $\gamma$ (default: $0.95$), stop if
$\PP(p < \alpha \mid k,b) \ge \gamma$ or $\PP(p \ge \alpha \mid k,b) \ge \gamma$.

\begin{algorithm}[t]
  \caption{Adaptive sequential permutation test (heuristic)}
  \label{alg:adaptive-sequential}
  \begin{algorithmic}[1]
    \Require Statistic $T$, observed value $T_0$, max resamples $B$, threshold $\alpha$, confidence $\gamma$
    \State $\texttt{min\_resamples} \gets \lceil 1/\alpha \rceil$; $B \gets \max(B,\texttt{min\_resamples})$; $k \gets 0$
    \For{$b = 1,\dots,B$}
      \State Sample a random permutation and compute $T_b$
      \If{$T_b \ge T_0$} \State $k \gets k+1$ \EndIf
      \If{$b \ge \texttt{min\_resamples}$}
        \State $q \gets \PP(p < \alpha \mid k,b)$ under $p \mid k,b \sim \mathrm{Beta}(1+k,1+b-k)$
        \If{$q \ge \gamma$ or $1-q \ge \gamma$} \State \Return $(k+1)/(b+1)$ \EndIf
      \EndIf
    \EndFor
    \State \Return $(k+1)/(B+1)$
  \end{algorithmic}
\end{algorithm}

\paragraph{Left-tail variant (Stage B).}
For split testing, replace the exceedance check by $T_b \le T_0$.

\subsection{Computational heuristics (scope note)}
\label{app:methods-heuristics}

\citrees includes several heuristics intended to reduce computation. Some of
these can change the algorithm's behavior and can break the clean ``fixed
hypothesis family'' assumptions used by the paper's fixed-node theory. For
paper-facing p-value guarantees, use fixed-$B$ testing and avoid adaptive
hypothesis-family modifications (in particular, disable feature muting).

\subsubsection{Feature muting}
\label{app:methods-feature-muting}

Feature muting removes features deemed uninformative at a node from future
consideration within the subtree. Because muting decisions depend on observed
labels via Stage~A p-values, it can violate the label-independence condition
(A0.3) needed for clean fixed-node calibration statements when treated across
multiple nodes.

\subsubsection{Feature scanning}
\label{app:methods-feature-scanning}

When early stopping is enabled, feature scanning tests ``promising'' features
first by sorting candidates using a cheap, non-permutation score. This can
substantially reduce computation when a strong signal exists. When early
stopping is disabled (fixed-$B$ mode), ordering does not affect the result.

\subsubsection{Threshold scanning}
\label{app:methods-threshold-scanning}

Analogous to feature scanning, threshold scanning orders candidate thresholds
using a cheap deterministic score (e.g., weighted impurity) so that early
stopping can terminate quickly if a good split exists. In fixed-$B$ mode, the
ordering is irrelevant.

\subsection{Multiple testing correction}
\label{app:methods-multipletesting}

\paragraph{Bonferroni (default).}
At a fixed node $t$, Stage~A tests $m_t=|F_t|$ features and compares
$\min_{j\in F_t} p_{t,j}$ to $\alpha_{\mathrm{sel}}/m_t$ when
\texttt{adjust\_alpha\_selector} is enabled. Stage~B analogously compares the
best split p-value to $\alpha_{\mathrm{split}}/|C_{t,j_t^\star}|$ when
\texttt{adjust\_alpha\_splitter} is enabled.

\paragraph{Max-$T$ / Westfall--Young (multi-selector mode).}
When combining multiple selector statistics, \citrees supports a max-statistic
construction: within each permutation, take the maximum across selectors for a
fixed feature (and then compute a single p-value from the max-statistics).
Under exchangeability, the resulting +1 p-value remains super-uniform (see
\cref{prop:multiselector-validity} and \cref{app:multiselector}). This is
closely related to the max-$T$ resampling principle \citep{WestfallYoung1993ResamplingBasedMultipleTesting}.

\subsection{Forests and resampling}
\label{app:methods-forests}

\citrees forests fit $M$ trees and aggregate their predictions. Each tree is a
conditional inference tree trained on either the full dataset or a resampled
dataset (bootstrap). Classic and Bayesian bootstrap are described in
\cref{sec:bootstrap}.

\begin{algorithm}[t]
  \caption{Forest training (bagging with optional bootstrap)}
  \label{alg:forest}
  \begin{algorithmic}[1]
    \Require Training data $(X,y)$ with $n$ rows, number of trees $M$, tree hyperparameters $\theta_{\mathrm{tree}}$
    \Require Optional bootstrap method, optional \texttt{max\_samples}
    \For{$m = 1,\dots,M$}
      \If{bootstrap enabled}
        \State Draw bootstrap indices $I_m$ by the chosen method (size $n$, with replacement)
        \If{\texttt{max\_samples} $< n$}
          \State Subsample $I_m$ without replacement down to \texttt{max\_samples}
        \EndIf
        \State Fit $T_m \gets \textsc{FitTree}(X[I_m], y[I_m]; \theta_{\mathrm{tree}})$
      \Else
        \State Fit $T_m \gets \textsc{FitTree}(X, y; \theta_{\mathrm{tree}})$
      \EndIf
    \EndFor
    \State \Return $\{T_1,\dots,T_M\}$
  \end{algorithmic}
\end{algorithm}

\paragraph{Prediction.}
Forest prediction is standard bagging: average class probability vectors for
classification and average scalar predictions for regression.
